\subsection*{ソフトウェアリリース管理と提供}
\addcontentsline{toc}{subsection}{6 ソフトウェアリリース管理と提供}
ここでいうリリースは、開発活動の外へソフトウェアと関連成果物を配布することである。社内リリースや顧客への配布を含む。複数のプラットフォーム版や機能差のある版がある場合、特定版の再現、正しい材料の包装、正しい対象への提供が必要になる。

\subsubsection*{ソフトウェアビルド}
\addcontentsline{toc}{subsubsection}{6.1 ソフトウェアビルド}
ソフトウェアビルドは、正しい版のソフトウェア構成品目と適切な構成データを使って、顧客やテストチームなどに提供するソフトウェアパッケージを作る活動である。ビルド手順は、正しい順序で手順を実施するために必要である。

以前のリリースを正確に再構築できることも重要である。そのためには、ソースコードだけでなく、コンパイラ、ビルドツール、スクリプト、設定、ビルド手順も SCM の管理下に置く必要がある。

継続的統合では、変更がソース管理リポジトリへコミットされると、ビルドが自動的に実行される。パイプラインは、取得、コンパイル、リンク、静的解析、単体テスト、コード網羅率測定、文書抽出などを行う。ビルドに含まれる成果物を示す SBOM は重要な構成管理出力である。

\subsubsection*{ソフトウェアリリース管理}
\addcontentsline{toc}{subsubsection}{6.2 ソフトウェアリリース管理}
ソフトウェアリリース管理は、実行プログラム、文書、リリースノート、構成データなど、製品要素を識別し、包装し、提供する活動である。いつリリースするかは、解決した問題の重大度、過去リリースの欠陥密度、顧客通知の必要性などによって判断する。

リリースの物理的内容を記録する情報は、版記述文書である。リリースノートには、新機能、既知の問題、必要なプラットフォーム条件を記す。パッケージには、インストールまたはアップグレード手順も含める。完全性を保証するために、デジタル署名などを使うこともある。

継続的デリバリでは、ビルド成果物をインストールパッケージへまとめ、検証環境または本番環境へ配備するためのパイプラインを設ける。

\par\smallskip
\begin{center}
\centering
\IfFileExists{assets/generated_figures/ch08/fig-ch08-11.png}{\includegraphics[width=0.96\linewidth]{assets/generated_figures/ch08/fig-ch08-11.png}}{\fbox{\parbox[c][48mm][c]{0.9\linewidth}{\centering 画像生成待ち\\ビルドからリリースまでのパイプライン図}}}
\par\smallskip
\refstepcounter{figure}\label{fig:ch08-11}図 \thefigure: ビルドからリリースまでのパイプライン図
\end{center}
図 \thefigure は、ビルドからリリースまでのパイプライン図を視覚的に整理したものである。
\par\smallskip
